% ============================================================
\chapter{Lois Usuelles et leurs Propri\'et\'es}
\label{ch:lois}
% ============================================================

La th\'eorie des probabilit\'es a produit, au fil des si\`ecles, une galerie de distributions remarquables --- chacune n\'ee d'un probl\`eme concret, et chacune dot\'ee de propri\'et\'es sp\'ecifiques qui la rendent irremplaçable dans son domaine. La loi normale de Gauss, la loi exponentielle des temps d'attente, la loi de Poisson des \'ev\'enements rares, la loi beta de la statistique bay\'esienne, la loi gamma de la th\'eorie des files d'attente~: ce chapitre les rassemble dans un panorama organis\'e, mettant en lumi\`ere leurs relations mutuelles (somme, conditionnement, passage \`a la limite) et les propri\'et\'es qui les caract\'erisent.

\begin{intuition}
Ce chapitre pr\'esente un catalogue organis\'e des lois de probabilit\'e les plus
courantes, tant discr\`etes que continues, avec leurs propri\'et\'es, relations
mutuelles et domaines d'application.
\end{intuition}

% ------------------------------------------------------------
\section{Panorama des lois discr\`etes}
% ------------------------------------------------------------

\subsection{Loi binomiale n\'egative}

\begin{definition}[Binomiale n\'egative]
$X \sim \text{BN}(r, p)$~: nombre d'essais pour obtenir le $r$-i\`eme succ\`es.
\[
    \PP(X = k) = \binom{k-1}{r-1} p^r (1-p)^{k-r}, \quad k = r, r+1, \ldots
\]
$\E[X] = r/p$, $\Var(X) = r(1-p)/p^2$. Pour $r = 1$, on retrouve la loi g\'eom\'etrique.
\end{definition}

\subsection{Loi uniforme discr\`ete}

\begin{definition}[Uniforme discr\`ete]
$X \sim \mathcal{U}\{a, a+1, \ldots, b\}$~:
\[
    \PP(X = k) = \frac{1}{b - a + 1}, \quad k = a, a+1, \ldots, b.
\]
$\E[X] = (a+b)/2$, $\Var(X) = ((b-a+1)^2 - 1)/12$.
\end{definition}

\subsection{Loi de Poisson --- compl\'ements}

\begin{proposition}[Stabilit\'e par somme]
Si $X_1 \sim \text{Poisson}(\lambda_1)$ et $X_2 \sim \text{Poisson}(\lambda_2)$
sont ind\'ependantes, alors $X_1 + X_2 \sim \text{Poisson}(\lambda_1 + \lambda_2)$.
\end{proposition}

\begin{proposition}[D\'ecomposition de Poisson]
Si $X \sim \text{Poisson}(\lambda)$ et chaque \'ev\'enement est class\'e de type~1
avec probabilit\'e $p$ et de type~2 avec probabilit\'e $1-p$, ind\'ependamment,
alors les nombres $N_1, N_2$ de chaque type sont ind\'ependants avec
$N_1 \sim \text{Poisson}(\lambda p)$ et $N_2 \sim \text{Poisson}(\lambda(1-p))$.
\end{proposition}

% ------------------------------------------------------------
\section{Panorama des lois continues}
% ------------------------------------------------------------

\subsection{Loi log-normale}

\begin{definition}[Log-normale]
$X \sim \text{LogN}(\mu, \sigma^2)$ si $\ln X \sim \mathcal{N}(\mu, \sigma^2)$.
\[
    f_X(x) = \frac{1}{x\sigma\sqrt{2\pi}}
    \exp\!\left(-\frac{(\ln x - \mu)^2}{2\sigma^2}\right), \quad x > 0.
\]
$\E[X] = e^{\mu + \sigma^2/2}$, $\Var(X) = e^{2\mu+\sigma^2}(e^{\sigma^2}-1)$.
\end{definition}

\subsection{Loi de Cauchy}

\begin{definition}[Cauchy]
$X \sim \text{Cauchy}(x_0, \gamma)$~:
\[
    f_X(x) = \frac{1}{\pi\gamma\left(1 + \left(\frac{x-x_0}{\gamma}\right)^2\right)}.
\]
L'esp\'erance et la variance n'existent pas.
\end{definition}

\begin{attention}[title=\textbf{Cauchy n'a pas de moments}]
La loi de Cauchy montre que toute variable al\'eatoire n'a pas forc\'ement
d'esp\'erance. La moyenne empirique de variables de Cauchy i.i.d.\ ne converge
pas~: la loi des grands nombres ne s'applique pas.
\end{attention}

\subsection{Loi de Weibull}

\begin{definition}[Weibull]
$X \sim \text{Weibull}(k, \lambda)$ avec $k, \lambda > 0$~:
\[
    f_X(x) = \frac{k}{\lambda}\left(\frac{x}{\lambda}\right)^{k-1}
    e^{-(x/\lambda)^k}\, \mathbf{1}_{(0,+\infty)}(x).
\]
Pour $k = 1$, on retrouve la loi exponentielle de param\`etre $1/\lambda$.
\end{definition}

\subsection{Loi chi-deux}

\begin{definition}[Chi-deux]
Si $Z_1, \ldots, Z_n$ sont i.i.d.\ $\mathcal{N}(0,1)$, alors
$\chi^2_n = Z_1^2 + \cdots + Z_n^2 \sim \Gamma(n/2, 1/2)$ est la loi du
\textbf{chi-deux \`a $n$ degr\'es de libert\'e}. $\E[\chi^2_n] = n$,
$\Var(\chi^2_n) = 2n$.
\end{definition}

\subsection{Loi de Student}

\begin{definition}[Student]
Si $Z \sim \mathcal{N}(0,1)$ et $V \sim \chi^2_n$ sont ind\'ependants, alors
$T = Z/\sqrt{V/n}$ suit la \textbf{loi de Student \`a $n$ degr\'es de libert\'e}.
$\E[T] = 0$ (pour $n > 1$), $\Var(T) = n/(n-2)$ (pour $n > 2$).
Pour $n = 1$, on retrouve la loi de Cauchy standard.
\end{definition}

\subsection{Loi de Fisher}

\begin{definition}[Fisher]
Si $U \sim \chi^2_m$ et $V \sim \chi^2_n$ sont ind\'ependants, alors
$F = (U/m)/(V/n)$ suit la \textbf{loi de Fisher \`a $(m, n)$ degr\'es de libert\'e}.
\end{definition}

% ------------------------------------------------------------
\section{Relations entre les lois}
% ------------------------------------------------------------

\begin{formules}[title=\textbf{Carte des relations entre lois}]
\begin{itemize}
    \item $\text{Bernoulli}(p) = \text{Bin}(1, p)$.
    \item $\text{Geom}(p) = \text{BN}(1, p)$.
    \item $\text{Exp}(\lambda) = \Gamma(1, \lambda)$.
    \item $\chi^2(n) = \Gamma(n/2, 1/2)$.
    \item $\text{Bin}(n, p) \xrightarrow{n\to\infty,\, np\to\lambda} \text{Poisson}(\lambda)$.
    \item $\text{Bin}(n, p) \xrightarrow{n\to\infty} \mathcal{N}(np, np(1-p))$ \quad (TCL).
    \item $\text{Poisson}(\lambda) \xrightarrow{\lambda\to\infty} \mathcal{N}(\lambda, \lambda)$.
    \item $t_n \xrightarrow{n\to\infty} \mathcal{N}(0,1)$.
    \item $\Gamma(n, \lambda)$ = loi de la somme de $n$ v.a.\ $\text{Exp}(\lambda)$ ind\'ependantes.
\end{itemize}
\end{formules}

% ------------------------------------------------------------
\section{Tableau r\'ecapitulatif g\'en\'eral}
% ------------------------------------------------------------

\begin{formules}[title=\textbf{R\'ecapitulatif des lois continues}]
\begin{center}
\renewcommand{\arraystretch}{1.4}
\begin{tabular}{lcccc}
    \toprule
    \textbf{Loi} & \textbf{Densit\'e} & \textbf{Support} & $\E[X]$ & $\Var(X)$ \\
    \midrule
    $\mathcal{U}([a,b])$ & $\frac{1}{b-a}$ & $[a,b]$ & $\frac{a+b}{2}$ & $\frac{(b-a)^2}{12}$ \\[4pt]
    $\text{Exp}(\lambda)$ & $\lambda e^{-\lambda x}$ & $[0,\infty)$ & $\frac{1}{\lambda}$ & $\frac{1}{\lambda^2}$ \\[4pt]
    $\mathcal{N}(\mu,\sigma^2)$ & $\frac{1}{\sigma\sqrt{2\pi}}e^{-(x-\mu)^2/(2\sigma^2)}$ & $\R$ & $\mu$ & $\sigma^2$ \\[4pt]
    $\Gamma(\alpha,\lambda)$ & $\frac{\lambda^\alpha}{\Gamma(\alpha)}x^{\alpha-1}e^{-\lambda x}$ & $(0,\infty)$ & $\frac{\alpha}{\lambda}$ & $\frac{\alpha}{\lambda^2}$ \\[4pt]
    $\text{Beta}(\alpha,\beta)$ & $\frac{x^{\alpha-1}(1-x)^{\beta-1}}{B(\alpha,\beta)}$ & $(0,1)$ & $\frac{\alpha}{\alpha+\beta}$ & $\frac{\alpha\beta}{(\alpha+\beta)^2(\alpha+\beta+1)}$ \\
    \bottomrule
\end{tabular}
\end{center}
\end{formules}

% ------------------------------------------------------------
\section{M\'ethodes de reconnaissance}
% ------------------------------------------------------------

\begin{remarque}
Pour identifier une loi, on peut proc\'eder ainsi~:
\begin{enumerate}
    \item Identifier le \textbf{support}~: $\N$, $\{0,1,\ldots,n\}$, $(0,\infty)$,
          $\R$, $(0,1)$, etc.
    \item Rep\'erer la \textbf{forme fonctionnelle}~: exponentielle d\'ecroissante,
          cloche sym\'etrique, puissance\ldots
    \item V\'erifier les \textbf{param\`etres} en calculant esp\'erance et variance.
\end{enumerate}
\end{remarque}

% ------------------------------------------------------------
\section{Exercices}
% ------------------------------------------------------------

\begin{exercice}
Soit $X \sim \Gamma(n, \lambda)$ avec $n \in \N^*$. Montrer que $X$ a m\^eme loi
que la somme de $n$ variables i.i.d.\ $\text{Exp}(\lambda)$.
\end{exercice}

\begin{exercice}
Montrer que si $X \sim \mathcal{N}(0,1)$, alors $X^2 \sim \chi^2(1)$.
\end{exercice}

\begin{exercice}
Soit $X \sim \text{Poisson}(\lambda)$ avec $\lambda = 100$. En utilisant
l'approximation normale, estimer $\PP(90 \leq X \leq 110)$.
\end{exercice}

\begin{exercice}
Soit $U \sim \mathcal{U}([0,1])$. Calculer la loi de $X = -\ln U$.
\end{exercice}

\begin{exercice}
Soit $X \sim \text{Beta}(1,1)$. Montrer que $X \sim \mathcal{U}([0,1])$.
\end{exercice}

\begin{exercice}
On observe le nombre de clients arrivant dans un magasin chaque heure, suppos\'e
suivre une loi de Poisson de param\`etre $\lambda = 12$. Chaque client ach\`ete
ind\'ependamment avec probabilit\'e $p = 0{,}3$. Quelle est la loi du nombre
de clients acheteurs par heure~?
\end{exercice}

\begin{exercice}
Montrer que si $X \sim t_n$ (Student \`a $n$ ddl), alors $X^2 \sim F(1, n)$
(Fisher).
\end{exercice}

\begin{exercice}
Soit $X$ de loi log-normale $\text{LogN}(0, 1)$. Calculer $\E[X]$, $\Var(X)$
et la m\'ediane de $X$.
\end{exercice}
